% --- PAGE 14: HALF-TITLE INTRODUCTION ---
\clearpage
\pagenumbering{arabic}
\setcounter{page}{1}
\pagestyle{fancy}
\fancyhf{}
\renewcommand{\headrulewidth}{0pt}
\cfoot{\fontsize{12pt}{14pt}\selectfont\thepage}

\vspace*{\fill}
\begin{center}
  {\fontsize{18pt}{24pt}\selectfont\bfseries INTRODUCTION\par}
\end{center}
\vspace*{\fill}
\newpage

% --- CHAPTER 1 INTRODUCTION ---
\begin{center}
  {\fontsize{16pt}{20pt}\selectfont\bfseries CHAPTER 1\par}
  \vspace{0.3cm}
  {\fontsize{16pt}{20pt}\selectfont\bfseries INTRODUCTION\par}
\end{center}
\vspace{0.4cm}

{\fontsize{14pt}{18pt}\selectfont\bfseries 1.1 BACKGROUND AND DOMAIN SIGNIFICANCE\par}
\vspace{0.3cm}

{\fontsize{12pt}{17.8pt}\selectfont
In modern engineering and technology institutions, the effective governance of continuous academic evaluations, laboratory deliverables, and extracurricular algorithmic preparation forms the foundation of student engineering outcomes. Computer science, information technology, and software engineering curricula demand that students actively engage not only with traditional theoretical coursework, but also with rigorous hands-on programming labs, full-stack mini-projects, competitive algorithmic problem solving, and industry-sponsored technical challenges.

However, existing operational methodologies within academic departments remain critically fragmented and error-prone. Departmental faculties, class advisors, and laboratory coordinators frequently depend on disparate, uncoordinated communication channels---including informal WhatsApp broadcast groups, Google Classrooms, shared cloud spreadsheets, and physical paper-based observation notebooks---to assign, collect, and evaluate student submissions. This disconnected approach generates severe operational bottlenecks: student submissions are frequently lost or misattributed, multi-week grading backlogs accumulate, and student accountability degrades due to the total absence of a centralized, tamper-evident audit trail.

Simultaneously, the global software recruitment ecosystem has shifted toward competency-based hiring paradigms where candidate evaluation is strictly driven by verified coding proficiency, active problem-solving streaks on platforms such as LeetCode, and continuous version-controlled contribution histories on GitHub. In conventional college setups, academic mentors and placement officers operate in an informational blind spot regarding students' daily algorithmic practice, unable to detect momentum loss or identify students in need of academic intervention. To bridge these operational and pedagogical gaps, this project introduces \textbf{VSBEC IT Vault}, an integrated enterprise platform engineered to enforce academic task governance, automate competitive telemetry harvesting, and host proctored sandboxed coding evaluations for institutional recruitment.

\vspace{0.4cm}
{\fontsize{14pt}{18pt}\selectfont\bfseries 1.2 TECHNICAL MOTIVATION AND PROBLEM STATEMENT\par}
\vspace{0.3cm}

Traditional Learning Management Systems (LMS) such as Moodle, Canvas, and Google Classroom were architected for general-purpose course content delivery rather than technical software engineering education. When applied to programming-heavy academic domains, these systems exhibit severe structural vulnerabilities:
\begin{enumerate}[noitemsep,topsep=3pt,leftmargin=*]
  \item \textbf{Absence of Native Compilation:} Conventional systems treat source code files as passive binary attachments. Evaluators must manually download, configure dependencies, and compile student programs on local workstations, introducing massive evaluation latency.
  \item \textbf{Single-Tier Insecure Authorization:} Most academic portals implement flat permission structures, lacking the granular Role-Based Access Control (RBAC) required to govern peer-coordinator reviews, faculty grading, and departmental HOD audits.
  \item \textbf{Vulnerability to Data Loss:} Academic records and submissions stored in ephemeral cloud containers or unbacked relational instances lack automated, offsite disaster recovery pipelines.
  \item \textbf{Disconnection from Industry Telemetry:} Existing tools provide zero automated hooks into competitive programming ecosystems, leaving institutions detached from corporate recruitment benchmarks.
\end{enumerate}

\vspace{0.4cm}
{\fontsize{14pt}{18pt}\selectfont\bfseries 1.3 PROJECT OBJECTIVES\par}
\vspace{0.3cm}

To resolve these technical and operational deficiencies, the \textbf{VSBEC IT Vault} project defines six concrete engineering objectives:
\begin{enumerate}[noitemsep,topsep=2pt,leftmargin=*]
  \item \textbf{Centralized Academic Task Governance:} Establish an auditable 3-tier review hierarchy (Student Coordinator $\rightarrow$ Faculty Class Advisor $\rightarrow$ Head of Department) with Cloudinary asset storage and automated Excel report compilers.
  \item \textbf{Sandboxed Multi-Language Compiler Engine:} Engineer an isolated execution jail supporting C, C++, Java 17, and Python 3 with sub-second execution timeout traps (3000ms), output buffer bounds (64 KB), and Judge0 cloud failover.
  \item \textbf{Automated Telemetry Harvesting Daemons:} Build background schedulers that poll LeetCode GraphQL and GitHub REST APIs daily to compute problem solves, active streaks, and class leaderboards.
  \item \textbf{Automated Disaster Recovery \& Dual-Schedule Backup Pipeline:} Deploy GitHub Actions workflows executing twice daily (07:50 AM IST and 11:55 PM IST) to export daily CSV/JSON progress reports, generate database snapshots, and enforce dual-remote Git synchronization.
  \item \textbf{High-Availability Multi-Account Brevo Notification Engine:} Deploy an intelligent multi-node Brevo email engine with automated failover for 6-digit OTP verification, alongside Telegram bot webhooks and VAPID Web Push notifications.
  \item \textbf{Non-Invasive Anti-Cheat Proctoring:} Integrate picture-in-picture (PIP) webcam monitoring, HTML5 fullscreen lockouts, and tab-switch telemetry event traps for recruitment evaluations.
\end{enumerate}

\vspace{0.4cm}
{\fontsize{14pt}{18pt}\selectfont\bfseries 1.4 ORGANIZATION OF THE REPORT\par}
\vspace{0.3cm}

The remainder of this report is structured as follows: Chapter 2 reviews relevant academic and industry literature. Chapter 3 analyzes existing system flaws and formalizes the proposed architecture and comparison matrix. Chapter 4 specifies hardware and software requirements. Chapter 5 details the core algorithmic, database, and telemetry implementation. Chapter 6 presents UML and DFD system designs. Chapter 7 validates test cases and testing methodologies. Chapter 8 evaluates feasibility and engineering methodology. Chapter 9 concludes the study with five future engineering extensions.\par}
\newpage
